% UBT-AI-PROVENANCE-BEGIN
% schema: ubt-ai-provenance/v1
% tier: C_working
% ai_assistance: disclosed
% human_review: risk-based
% editorial_responsibility: Ing. David Jaroš
% policy: ../../../AI_PROVENANCE.md
% notice: Working material; exhaustive human review is not claimed.
% UBT-AI-PROVENANCE-END

\chapter{Theta, Mellin, zeta, cesty a prvočíselná struktura}
\label{ch:theta-mellin-feynman}

\section{Proč tato kapitola existuje}

Tato kapitola je záměrně širší než kanonické dokumenty tvrzení UBT. Jejím
pedagogickým účelem je umístit několik klasických matematických struktur vedle
redukovaného propojení UBT s theta funkcemi, aby byly jejich vztahy patrné na
jednom místě. Výsledek se \emph{nestává} teorémem UBT jen proto, že je v UBT
užitečný.

V próze proto používáme tři stavové štítky:
\begin{itemize}
  \item\textbf{Klasické}: standardní matematika nebo standardní kvantová mechanika;
  \item\textbf{Propojení}: přesné tvrzení o redukovaném modelu theta, který
        propojuje standardní matematiku se zápisem UBT;
  \item\textbf{Otevřeno v UBT}: tvrzení, které by muselo být odvozeno z kanonického
        pole, rovnic nebo akce UBT, než by se mohlo stát fyzikálním tvrzením UBT.
\end{itemize}

Ústředním redukovaným objektem je
\begin{equation}
  S(t,\psi)
  = \sum_{n\ge 1} a_n\,e^{-\pi\psi n^2}e^{i\pi t n^2}
  = \sum_{n\ge1}a_n e^{-\pi(\psi-it)n^2}.
  \label{eq:reduced-theta-amplitude}
\end{equation}
Jde o stejné jádro s redukovaným komplexním časem, jaké se používá v
\path{canonical/bridges/theta_complex_time_bridge.tex}. V této kapitole jde o
\textbf{model propojení}; nepředpokládá se, že kanonické hlavní pole UBT
$\Theta(q,\tau)$ je Jacobiho theta funkcí.

\paragraph{Znaménková konvence použitá v souvisejících poznámkách k RH.}
Některé dřívější pracovní poznámky UBT/RH používaly kladnou tepelnou proměnnou
$y=-\psi$, a proto zapisovaly parametr theta jako $t-i\psi=t+iy$ s konvergencí
na odpovídající polopřímce $y>0$. Tato kapitola zapisuje tlumicí proměnnou přímo
jako $\psi>0$ v rovnici~\eqref{eq:reduced-theta-amplitude}. Transformační
identity jsou po přeznačení $y=-\psi$ stejné; znaménkovou konvenci je však třeba
uvést dříve, než se k některé polopřímce připojí fyzikální výklad UBT.

\section{Jedno spektrum, několik reprezentací}

\subsection{Ortogonalita módů versus neortogonalita theta šablon}

Pro pevné $\psi$ jsou základní časové módy v rovnici~\eqref{eq:reduced-theta-amplitude}
\begin{equation}
  e^{i\pi n^2 t}.
\end{equation}
Na intervalu $t\in[0,2]$ jsou ortogonální:
\begin{equation}
 \int_0^2 e^{i\pi(n^2-m^2)t}\,dt
 = 2\,\delta_{nm},
 \qquad n,m\ge1.
 \label{eq:quadratic-mode-orthogonality}
\end{equation}
Důležitý rozdíl spočívá v tom, že \emph{úplné šablony}
$S_\psi(t):=S(t,\psi)$ při různých hodnotách $\psi$ obecně nejsou ortogonální.

\begin{theorem}[Klasická identita Gramova jádra]
Pro vnitřní produkt
\begin{equation}
  \langle f,g\rangle:=\frac12\int_0^2 f(t)\overline{g(t)}\,dt,
\end{equation}
a amplitudy $S_\psi$ z rovnice~\eqref{eq:reduced-theta-amplitude} platí
\begin{equation}
  G(\psi,\phi)
  :=\langle S_\psi,S_\phi\rangle
  =\sum_{n\ge1}|a_n|^2e^{-\pi(\psi+\phi)n^2}.
  \label{eq:gram-general}
\end{equation}
Pro $a_n\equiv1$,
\begin{equation}
  G(\psi,\phi)
  =\frac{\vartheta_3(0\mid i(\psi+\phi))-1}{2}.
  \label{eq:gram-theta}
\end{equation}
\end{theorem}

\begin{proof}
Dosazením obou součtů do vnitřního součinu rovnice
\eqref{eq:quadratic-mode-orthogonality} odstraní všechny smíšené členy
$m\neq n$ a ponechá
\[
  \sum_n |a_n|^2 e^{-\pi\psi n^2}e^{-\pi\phi n^2}.
\]
U jednotkových vah standardní identita
$\vartheta_3(0\mid iu)=1+2\sum_{n\ge1}e^{-\pi u n^2}$ dává
rovnici~\eqref{eq:gram-theta}.
\end{proof}

\paragraph{Praktický důsledek.}
Neortogonalita šablon není překážkou detekce.  Pokud je pozorovaný signál
$f(t)$, normalizované skóre shodného filtru je
\begin{equation}
  \rho(\psi,\Delta)
  =\frac{\left|\int f(t)\overline{S(t-\Delta,\psi)}\,dt\right|^2}
  {\|f\|^2\,\|S_\psi\|^2}.
  \label{eq:theta-matched-filter}
\end{equation}
Nejlepší samostatná šablona se získá maximalizací $\rho$ vzhledem k $(\psi,\Delta)$.
Ortogonalita je důležitá při současném prokládání mnoha šablon; tehdy
Gramova matice musí být invertována nebo regularizována, například pomocí $G+\lambda I$.

\subsection{Fourierova transformace v reálném čase}

Zvolme Fourierovu konvenci
\begin{equation}
  \widehat f(\omega)=\int_{-\infty}^{\infty}f(t)e^{-i\omega t}\,dt.
\end{equation}
Formálně nebo distribučně,
\begin{equation}
  \widehat S(\omega,\psi)
  =2\pi\sum_{n\ge1}a_n e^{-\pi\psi n^2}
       \delta(\omega-\pi n^2).
  \label{eq:fourier-delta-spectrum}
\end{equation}
Spektrum kvadratických módů je tedy již Diracovým hřebenem ve frekvenci
reálného času; k vytvoření delta čar není třeba demodulovat jednotlivé módy.

\subsection{Laplaceova transformace v imaginárním čase}

Pro $\Re z$ dostatečně velký,
\begin{align}
  \mathcal L_\psi[S](z;t)
   &=\int_0^\infty e^{-z\psi}S(t,\psi)\,d\psi \\
   &=\sum_{n\ge1}
      \frac{a_n e^{i\pi t n^2}}{z+\pi n^2}.
  \label{eq:laplace-resolvent}
\end{align}
Stejná vlastní čísla, která vytvořila Fourierovy delta čáry, se nyní objevují jako póly rezolventy
$z=-\pi n^2$.

\subsection{Mellinova transformace v imaginárním čase}

Pro $\Re s$ v polorovině absolutní konvergence platí
\begin{align}
  \mathcal M_\psi[S](s;t)
  &:=\int_0^\infty \psi^{s-1}S(t,\psi)\,d\psi \\
  &=\Gamma(s)\pi^{-s}
    \sum_{n\ge1}\frac{a_n e^{i\pi t n^2}}{n^{2s}}.
  \label{eq:mellin-time-twisted}
\end{align}
Identita vyplývá z
\begin{equation}
 \int_0^\infty \psi^{s-1}e^{-\pi n^2\psi}\,d\psi
 =\Gamma(s)(\pi n^2)^{-s}.
\end{equation}

Pro jádro Gram s $a_n\equiv1$ a $u=\psi+\phi$,
\begin{equation}
 \int_0^\infty u^{s-1}G(u)\,du
 =\Gamma(s)\pi^{-s}\zeta(2s).
 \label{eq:gram-mellin-zeta}
\end{equation}
Jde o klasickou matematiku theta funkcí a Mellinovy transformace. Její význam
je zde organizační: tatáž redukovaná banka theta funkcí, která vytváří
neortogonální šablony, je také tepelnou stopou, jejíž Mellinovou reprezentací je
spektrální zeta funkce.

\paragraph{Spektrální trojúhelník.}
Rovnice~\eqref{eq:fourier-delta-spectrum},\eqref{eq:laplace-resolvent} a
\eqref{eq:mellin-time-twisted} poskytují tři pohledy na stejné kvadratické spektrum:
\begin{equation}
 \boxed{
 \begin{array}{ccl}
   \mathcal F_t   &\longrightarrow& \delta(\omega-\pi n^2),\\[2mm]
   \mathcal L_\psi&\longrightarrow& (z+\pi n^2)^{-1},\\[2mm]
   \mathcal M_\psi&\longrightarrow& \Gamma(s)\pi^{-s}n^{-2s}.
 \end{array}}
 \label{eq:spectral-triangle}
\end{equation}

\section{Mellin je Fourierova transformace v logaritmické souřadnici}

Nechť
\begin{equation}
  M_f(s)=\int_0^\infty f(x)x^{s-1}\,dx,
  \qquad s=\sigma+i\omega.
\end{equation}
Položme $x=e^u$, tedy $dx=e^u du$. Pak
\begin{equation}
  M_f(\sigma+i\omega)
  =\int_{-\infty}^{\infty}
   f(e^u)e^{\sigma u}e^{i\omega u}\,du.
  \label{eq:mellin-log-fourier}
\end{equation}
Kromě konvenčního znaku Fourierovy fáze se jedná o Fourierovu transformaci
$u=\ln x$ váženého signálu $f(e^u)e^{\sigma u}$.  Mellinova analýza je tedy
přirozená harmonická analýza pro multiplikativní škálování, zatímco běžná Fourierova analýza
je přirozená pro aditivní posun.

Toto pozorování je zvlášť užitečné při interpretaci zeta funkcí: svislá
souřadnice $\Im s$ v komplexní rovině je logaritmickou frekvencí.

\section{Současný řez a časově fázovaná rodina zeta funkcí}
\label{sec:time-twisted-zeta}

\paragraph{Souřadnicová konvence UBT.}
V pracovním výkladu komplexního času použitém v tomto výzkumném směru se řez
současnosti volí jako $t=0$. Jde o souřadnicovou konvenci UBT, nikoli o tvrzení,
že Riemannova zeta funkce vnitřně rozpoznává fyzikální „přítomnost“.

Pro jednotkové váhy $a_n\equiv1$ definujme kvadratické časové zkroucení
\begin{equation}
  Z_\Theta(s;t)
  :=\sum_{n=1}^{\infty}
    \frac{e^{i\pi t n^2}}{n^{2s}}.
  \label{eq:time-twisted-zeta}
\end{equation}
Rovnice~\eqref{eq:mellin-time-twisted} pak přejde do tvaru
\begin{equation}
 \boxed{
  \mathcal M_\psi[S](s;t)
  =\Gamma(s)\pi^{-s}Z_\Theta(s;t).}
 \label{eq:mellin-Ztheta}
\end{equation}
Na řezu současnosti platí
\begin{equation}
 \boxed{Z_\Theta(s;0)=\zeta(2s).}
 \label{eq:present-zeta}
\end{equation}
Obyčejná zeta je tedy jedním členem --- členem s $t=0$ --- přirozené rodiny
vytvořené týmž jádrem theta s redukovaným komplexním časem.

\begin{remark}
Exponenciální faktor v rovnici~\eqref{eq:time-twisted-zeta} je \emph{aditivní
kvadratické zkroucení}. Pro obecné $t$ nemá výsledná Dirichletova řada obvyklý
Eulerův součin. Dynamiku prvočíselných faktorů nelze odvozovat pouze z existence
Eulerova součinu při $t=0$.
\end{remark}

\subsection{Jednoduché speciální časy}

Rodina je pod $t\mapsto t+2$ periodická, protože $n^2\in\mathbb Z$:
\begin{equation}
 Z_\Theta(s;t+2)=Z_\Theta(s;t).
\end{equation}
V $t=1$,
\begin{align}
 Z_\Theta(s;1)
 &=\sum_{n\ge1}\frac{(-1)^{n^2}}{n^{2s}}
 =\sum_{n\ge1}\frac{(-1)^n}{n^{2s}}\\
 &=-\eta(2s)
 =-\bigl(1-2^{1-2s}\bigr)\zeta(2s),
 \label{eq:t1-eta}
\end{align}
kde $\eta$ je Dirichletova eta funkce.

\subsection{Racionální časy: rozklad pomocí Hurwitzovy zeta funkce}

Nechť $t=a/q$, kde $a$ a $q>0$ jsou celá čísla. Koeficient
\begin{equation}
 c_n=e^{i\pi a n^2/q}
\end{equation}
je periodický s periodou $M=2q$ (skutečná perioda může být někdy menší).
Rozdělením Dirichletovy řady do zbytkových tříd dostaneme přesnou klasickou
identitu
\begin{equation}
 \boxed{
 Z_\Theta\!\left(s;\frac aq\right)
 =M^{-2s}\sum_{r=1}^{M}
 e^{i\pi a r^2/q}
 \zeta\!\left(2s,\frac rM\right),
 \qquad M=2q,}
 \label{eq:rational-hurwitz}
\end{equation}
kde $\zeta(s,\alpha)$ je Hurwitzova zeta funkce.

\begin{proof}
Každé $n\ge1$ lze jednoznačně zapsat jako $n=mM+r$ s $m\ge0$ a $1\le r\le M$.
Periodicita dává $c_{mM+r}=c_r$ a
\begin{align*}
 Z_\Theta(s;a/q)
 &=\sum_{r=1}^{M}c_r\sum_{m=0}^{\infty}(mM+r)^{-2s}\\
 &=M^{-2s}\sum_{r=1}^{M}c_r
   \zeta\!\left(2s,\frac rM\right).
\end{align*}
\end{proof}

Díky tomu jsou „alternativní funkce zeta“ generované mimo $t=0$ přesně určené.
Běžná zeta, eta, Hurwitzova zeta a kvadratické Gaussovy součty zde nejsou
nesouvisející objekty: jsou to různé řezy nebo reprezentace jedné rodiny
kvadratických fází.

\section{Racionální návraty a kvadratické Gaussovy součty}
\label{sec:gauss-revivals}

Pro racionální $t=a/q$ závisí fáze $e^{i\pi a n^2/q}$ pouze na konečné zbytkové
třídě. Odpovídající konečné kvadratické součty jsou Gaussovy součty. Obvyklá
normalizace je
\begin{equation}
 g(a,q)=\sum_{r=0}^{q-1}e^{2\pi i a r^2/q}.
 \label{eq:gauss-sum}
\end{equation}
Pro nesoudělné moduly $q_1,q_2$ dává čínská věta o zbytcích při normalizaci z
rovnice~\eqref{eq:gauss-sum}
\begin{equation}
 g(a,q_1q_2)=g(aq_2,q_1)\,g(aq_1,q_2),
 \qquad \gcd(q_1,q_2)=1.
 \label{eq:gauss-crt}
\end{equation}
Složené jmenovatele se tedy rozkládají na nesoudělné části a mocniny prvočísel
jsou nerozložitelnými aritmetickými stavebními prvky struktury návratů v
racionálních časech.

Jde o slibnější místo pro hledání \emph{dynamického} původu prvočíselných
sektorů UBT než pouhé pozorování Eulerova součinu $\zeta(2s)$, protože
aritmetiku Gaussových součtů nese přímo kvadratická fáze theta.

\section{Součet theta módů a Feynmanův součet vinutí}
\label{sec:theta-feynman}

Vztah mezi theta funkcí a drahami je přesný v kontrolovaném učebnicovém
příkladu volné částice o hmotnosti $m$, která se pohybuje po kružnici o obvodu
$L$.

Vlastní funkce hybnosti a jejich energie jsou
\begin{equation}
 \phi_n(x)=L^{-1/2}e^{2\pi i n x/L},
 \qquad
 E_n=\frac{\hbar^2}{2m}\left(\frac{2\pi n}{L}\right)^2.
\end{equation}
Proto je propagátor v reálném čase
\begin{equation}
 K(x,x';T)
 =\frac1L\sum_{n\in\mathbb Z}
 \exp\!\left(\frac{2\pi i n(x-x')}{L}
             -\frac{iE_nT}{\hbar}\right).
 \label{eq:circle-mode-propagator}
\end{equation}
Jde o Jacobiho theta součet přes módy. Poissonova sumační formule jej převádí na
\begin{equation}
 \boxed{
 K(x,x';T)
 =\sqrt{\frac{m}{2\pi i\hbar T}}
  \sum_{w\in\mathbb Z}
  \exp\!\left[
   \frac{i m(x-x'+wL)^2}{2\hbar T}
  \right].}
 \label{eq:circle-winding-propagator}
\end{equation}
Celé číslo $w$ označuje sektory vinutí. Exponent obsahuje klasickou akci volné
částice pro zdviženou dráhu z $x'$ do $x+wL$:
\begin{equation}
 S_w^{\rm cl}=\frac{m(x-x'+wL)^2}{2T}.
\end{equation}
Tentýž propagátor proto připouští dvě přesné reprezentace:
\begin{equation}
 \boxed{
 K_\Sigma
 \quad\Longleftrightarrow\quad
 K_{\mathbb Z}}
 \label{eq:mode-winding-duality}
\end{equation}
Symbol $K_\Sigma$ označuje spektrální součet módů a $K_{\mathbb Z}$ součet přes
sektory vinutí,
s Poissonovou--Jacobiho transformací jako propojením.

\paragraph{Důležité upřesnění.}
Klasickou trajektorii vybírá \emph{stacionární} akce, $\delta S=0$, nikoli
nutně její minimum. Více stacionárních drah je v kvantové mechanice běžných;
samo o sobě to neimplikuje mnohovesmírnou interpretaci.

\section{Kde prvočísla vstupují a kde nevstupují}

\subsection{Eulerův součin na nefázovaném řezu}

Pro $t=0$ a jednotkové váhy platí
\begin{equation}
 Z_\Theta(s;0)=\zeta(2s)
 =\prod_{p}\left(1-p^{-2s}\right)^{-1},
 \qquad \Re s>\frac12.
 \label{eq:zeta-euler-product}
\end{equation}
Jde o hlubokou změnu reprezentace: od aditivního součtu přes všechna kladná
celá čísla k multiplikativnímu součinu přes prvočísla. Je to však klasická
vlastnost celých čísel, nikoli sama o sobě důkaz, že dynamika UBT vybírá
prvočísla. Pro obecné váhy $a_n$ má Dirichletova řada
$\sum |a_n|^2n^{-2s}$ Eulerův součin pouze tehdy, když koeficienty splňují
vhodné podmínky multiplikativity.

\subsection{Mřížky pólů lokálních Eulerových faktorů}

Jediný formální místní faktor
\begin{equation}
 L_p(s)=\left(1-p^{-2s}\right)^{-1}
\end{equation}
má póly, když $p^{-2s}=1$.  Zápis $s=\sigma+i\omega$ dává
\begin{equation}
 \sigma=0,
 \qquad
 \omega_{p,k}=\frac{\pi k}{\ln p},
 \qquad k\in\mathbb Z.
 \label{eq:local-prime-poles}
\end{equation}
Pro různá prvočísla $p\neq q$ se nenulové pólové frekvence nemohou shodovat. Je-li
$k/\ln p=\ell/\ln q$ s nenulovými celými čísly $k,\ell$, pak
$q^k=p^\ell$, což odporuje jednoznačnosti rozkladu na prvočísla. Všechny
lokální faktory sdílejí triviální polohu $k=0$.

Jde o póly \emph{jednotlivých formálních Eulerových faktorů}. Leží mimo
polorovinu konvergence Eulerova součinu a nesmějí se zaměňovat s póly
analytického pokračování globální funkce $\zeta(2s)$, jejíž pól leží v $s=1/2$.

Počítá-li $N(T;P)$ kladné pólové frekvence lokálních faktorů do výšky $T$ pro
prvočísla $p\le P$, pak
\begin{align}
 N(T;P)
 &=\sum_{p\le P}\left\lfloor\frac{T\ln p}{\pi}\right\rfloor\\
 &=\frac{T}{\pi}\vartheta(P)+O(\pi(P)),
 \label{eq:local-pole-count}
\end{align}
kde
\begin{equation}
 \vartheta(P)=\sum_{p\le P}\ln p
\end{equation}
je Čebyševova theta funkce. Z věty o prvočíslech plyne
$\vartheta(P)\sim P$.

\subsection{Vertikální Fourierovo spektrum a mocniny prvočísel}

Při pevném $\sigma>1/2$,
\begin{equation}
 \zeta\bigl(2(\sigma+i\omega)\bigr)
 =\sum_{n\ge1}n^{-2\sigma}e^{-i(2\ln n)\omega}.
 \label{eq:zeta-vertical-fourier}
\end{equation}
Fourierova transformace v $\omega$ tedy vytváří spektrální čáry v bodech
\begin{equation}
 \xi_n=2\ln n.
\end{equation}
Pro izolaci prvočísel a jejich mocnin použijeme logaritmickou derivaci:
\begin{equation}
 -\frac{\zeta'(z)}{\zeta(z)}
 =\sum_{n\ge1}\frac{\Lambda(n)}{n^z}
 =\sum_p\sum_{k\ge1}(\ln p)p^{-kz},
 \qquad \Re z>1.
 \label{eq:log-derivative-von-mangoldt}
\end{equation}
S $z=2(\sigma+i\omega)$,
\begin{equation}
 -\frac{\zeta'\!\left(2(\sigma+i\omega)\right)}
        {\zeta\!\left(2(\sigma+i\omega)\right)}
 =\sum_p\sum_{k\ge1}
   (\ln p)p^{-2k\sigma}
   e^{-i(2k\ln p)\omega}.
 \label{eq:prime-power-fourier-comb}
\end{equation}
Hřeben mocnin prvočísel tedy leží na frekvencích
\begin{equation}
 \boxed{\xi_{p,k}=2k\ln p.}
 \label{eq:prime-power-frequencies}
\end{equation}
Různé čáry se shodují pouze tehdy, když jsou totožné i příslušné mocniny
prvočísel.

\subsection{Počítání prvočísel a frekvence nul zeta funkce}

Funkce počítající vážené mocniny prvočísel
\begin{equation}
 \psi_C(x)=\sum_{n\le x}\Lambda(n)
\end{equation}
souvisí klasickým explicitním vzorcem s netriviálními nulami
$\rho=\beta+i\gamma$ funkce $\zeta$. Schematicky platí
\begin{equation}
 \psi_C(x)=x-\sum_\rho\frac{x^\rho}{\rho}+\cdots.
 \label{eq:explicit-formula-schematic}
\end{equation}
Zápis $x=e^u$ změní každý nulový příspěvek na
\begin{equation}
 x^\rho=e^{\beta u}e^{i\gamma u}.
\end{equation}
Imaginární části $\gamma$ nul zeta funkce tedy působí jako oscilační frekvence
v logaritmické proměnné $u=\ln x$. Jde o klasické spektrální propojení
prvočísel a nul zeta funkce.

\section{Co to znamená pro prvočíselné sektory UBT}

Nesmí se spojovat dva odlišné mechanismy.

\paragraph{Cesta A: stopová/orbitální cesta.}
Analogie se stopovými formulemi naznačují, že by prvočísla měla odpovídat
primitivním periodickým objektům, jejichž opakované průchody generují mocniny
prvočísel. Pro obyčejnou volnou částici na $S^1$ však spektrum délek
primitivních orbit není $\{\ln p\}$. Proto
\begin{quote}
\textbf{GAP-THETA-PRIME-1:} odvodit ze skutečného operátoru UBT primitivní
orbitální či spektrální strukturu, jejíž přirozená primitivní data jsou
$\ln p$; jinak netvrdit, že Eulerův součin dynamicky odvozuje prvočíselné
sektory UBT.
\end{quote}

\paragraph{Cesta B: cesta racionálních návratů.}
Samotná kvadratická dynamika theta má v racionálních časech aritmetickou
strukturu. Gaussovy součty se faktorizují přes nesoudělné jmenovatele, takže
mocniny prvočísel jsou přirozenými stavebními prvky. To motivuje
\begin{quote}
\textbf{GAP-THETA-PRIME-2:} přepsat existující konstrukci prvočíselných sektorů
UBT $\alpha$ do proměnných návratů a vinutí a ověřit, zda její rozlišená
prvočísla vycházejí z racionální časové struktury Gaussových součtů, a nikoli
z dodatečně vybraného seznamu.
\end{quote}

Obě cesty jsou slučitelné jako směry výzkumu, dosud však není známo, zda
popisují tentýž mechanismus.

\section{Operátorová otázka za celým obrazem}

Rovnice~\eqref{eq:reduced-theta-amplitude} může být zapsána formálně zavedením
\begin{equation}
 H_0|n\rangle=\pi n^2|n\rangle.
\end{equation}
Evoluci módů pak generuje operátor
\begin{equation}
 U_0(t,\psi)=e^{-\psi H_0}e^{+itH_0}.
 \label{eq:formal-complex-propagator}
\end{equation}
Redukovaná amplituda theta se tím skutečně podobá propagátoru. Rozhodující
otázka UBT je však silnější:
\begin{quote}
\textbf{GAP-THETA-PROP:} odvodit operátor $H_\Theta$ z kanonické akce UBT nebo
rovnic pole tak, aby příslušné jádro či maticový prvek
$e^{-\psi H_\Theta}e^{\pm itH_\Theta}$ poskytoval zde použitou strukturu theta.
\end{quote}
Dokud se tato mezera neuzavře, jsou vztahy mezi módy a vinutím, Mellinovou
transformací a zeta funkcí i mezi návraty a prvočísly rigorózní matematikou
redukovaného modelu propojení, nikoli odvozením dynamiky UBT.

\section{Mapa velkého obrazu}

Užitečná syntéza je tedy
\begin{equation}
\boxed{
\begin{array}{ccccc}
 & \mathcal K_\Theta & & &\\
 & \displaystyle\sum_n a_n e^{-\pi\psi n^2}e^{i\pi t n^2} & & &\\[2mm]
 \swarrow & \downarrow & \searrow & &\\[-1mm]
 \mathcal F_t/\mathcal L_\psi & \mathcal M_\psi & \mathcal P & &\\
 \mathcal D & Z_\Theta(s;t) & \mathcal W & &\\
 & \downarrow\ t=0 & \downarrow\ t\in\mathbb Q & &\\
 & \zeta(2s) & \mathcal G_{\mathbb Q} & &\\
 & \downarrow & & &\\
 & \mathcal E_p & & &
\end{array}}
\end{equation}
V diagramu je $\mathcal K_\Theta$ redukovaným jádrem theta, $\mathcal P$
Poissonovou--Jacobiho transformací, $\mathcal D$ větví delta čar a rezolventy,
$\mathcal W$ součtem vinutí a drah, $\mathcal G_{\mathbb Q}$ Gaussovou
aritmetikou racionálních časů a $\mathcal E_p$ větví Eulerova součinu, mocnin
prvočísel a nul.
Klasické šipky jsou přesné. Výzkumným úkolem je určit, které z nich přežijí po
nahrazení redukovaného jádra kanonickým operátorem UBT a zda některá z obou
prvočíselných cest vysvětluje existující sektory programu UBT $\alpha$.
